\chapter{Brute Force}

Since the dawn of man, whenever encountered with a difficult problem, we have
decided ``eh, let's do it in the hardest way possible because I can't think of a
better solution.'' The history of brute forcing things goes all the way back to
our ancestors discovering fire,
\footnote{This is probably not true.} 
all the way to our current solution to climate change.
\footnote{Which is doing nothing and hoping
everything sorts itself out, which is really disappointing.}

The idea of brute forcing something is really as simple as it gets. We have $n$
number of literals, and therefore $2^n$ possible assignments, and we will check
them all until we have either found a solution, or until we've run out of
assignments to check. Because of how simple the idea is, it would just be
better to show code and explain that instead of explaining then showing code.

\begin{algorithm}[H]
\caption{Brute Force}\label{alg:3}
\begin{algorithmic}[1]
\Procedure{SolveBruteForce}{}

\For{\texttt{each possible assignment[1..numVars]}}
    \State $solved \gets True$

    \For{$Clause[1..n]$ in $Clauses$}
        \State $satisfied \gets False$

        \For{$i \gets 1$ up to $n$}
            \State $v \gets Clause[i]$
            \State $idx \gets \Call{abs}{v}$

            \If{$v > 0$}
                \State $satisfied \gets assignment[idx] \vee satisfied$
            \Else 
                \State $satisfied \gets \neg assignment[idx] \vee satisfied$
            \EndIf
        \EndFor
        \State $solved \gets satisfied \wedge solved$
    \EndFor

    \If{$solved = True$}
        \State $vars \gets assignment$
        \State return $True$
    \EndIf
\EndFor

\State return $False$

\EndProcedure
\end{algorithmic}
\end{algorithm}

Let's understand what's going on. Firstly, an \textit{assignment} in this case 
is an array of boolean values. Under the current assignment, we assume that it's
satisfactory (innocent until proven guilty) as shown on \textbf{line 3}. Next, 
we iterate through each clause in our database, and we assume that the clause is 
not satisfied. This is because or-ing a false value with a true value will get us 
a satisfied clause, but or-ing a true value with all false values will keep us 
at true. And if we decided to ``and'' everything together, it would have its own 
issues as well. 

On \textbf{line 6}, we iterate through the current clause. We grab its literals at 
index \textit{i}, and then we take the absolute value of that literal. This is because 
the absolute values of our literals ranges from $[1..numVars]$ and each assignment has 
exactly $numVars$ variables in it. So by taking the absolute value of a literal, we 
get to use it as an index into our \textit{assignment} variable.

On \textbf{line 9}, we check if the literal is greater than 0. If it is greater than 
0, we just ``or'' its assignment's value with whether or not the clause is currently 
satisfied. Otherwise, we take the negation of the assignment's value and ``or'' that 
with whether or not the clause is currently satisfied. This is because negative literals 
means ``use the opposite value of whatever I've been assigned with.'' Then we 
``and'' that clause-satisfied value with whether or not the equation has been solved.
We repeat this until there are no more clauses.

Then on \textbf{line 17}, we check if we've solved the equation. If we have, we 
store that assignment in our SAT solver and return $true$ to signify that we've found 
a solution. Otherwise, we move on to the next assignment. We do that until there are no
more assignments or we've found a satisfactory assignment. If we reach a point where 
we've looked through every assignment and none of them are satisfactory, we return 
$false$ on \textbf{line 22}.

Since it's trivially easy to see why the code works, we will move onto something
more interesting: optimizing brute force.

\section{Optimizing Brute Force}

I'll be completely honest with you, if you want to skip this section, go ahead.
It's more or less just something interesting enough we can do to make the
(arguably) worst algorithm\footnote{I do not know of any algorithm worse than
brute forcing.} better.

Since we're dealing with pseudocode, we're not going to consider the hardware
side of computers such as branch prediction, cache locality, or anything that
the hardware uses to execute our code (other than the CPU). We will consider
it from a pure Big-O perspective (excluding constant factors unless we're
choosing between two algorithms with the same computational complexity).

The first thing we will add is short circuiting to our clause evaluation. After
\textbf{line 15} where we finish OR-ing the ``clause satisfied'' boolean, we will add in
a new block:

\begin{algorithm}[H]
\caption{Brute Force}\label{alg:4}
\begin{algorithmic}[1]
\Procedure{SolveBruteForce}{}

\For{\texttt{each possible assignment[1..numVars]}}
    \State $solved \gets True$

    \For{$Clause[1..n]$ in $Clauses$}
        \State $satisfied \gets False$

        \For{$i \gets 1$ up to $n$}
            \State $v \gets Clause[i]$
            \State $idx \gets \Call{abs}{v}$

            \If{$v > 0$}
                \State $satisfied \gets assignment[idx] \vee satisfied$
            \Else 
                \State $satisfied \gets \neg assignment[idx] \vee satisfied$
            \EndIf

            \If{$satisfied = True$}
                \State \texttt{Move onto the next clause}
            \EndIf
        \EndFor
        \State $solved \gets satisfied \wedge solved$
    \EndFor

    \If{$solved = True$}
        \State $vars \gets assignment$
        \State return $True$
    \EndIf
\EndFor

\State return $False$
\EndProcedure
\end{algorithmic}
\end{algorithm}

This just says ``once this clause is satisfied, go evaluate the next one
immediately.''

We can then add another conditional after this loop, which will move onto the
next variable assignment the instant the equation is no longer satisfied:

\begin{algorithm}[H]
\caption{Brute Force}\label{alg:5}
\begin{algorithmic}[1]
\Procedure{SolveBruteForce}{}

\For{\texttt{each possible assignment[1..numVars]}}
    \State $solved \gets True$

    \For{$Clause[1..n]$ in $Clauses$}
        \State $satisfied \gets False$

        \For{$i \gets 1$ up to $n$}
            \State $v \gets Clause[i]$
            \State $idx \gets \Call{abs}{v}$

            \If{$v > 0$}
                \State $satisfied \gets assignment[idx] \vee satisfied$
            \Else 
                \State $satisfied \gets \neg assignment[idx] \vee satisfied$
            \EndIf

            \If{$satisfied = True$}
                \State \texttt{Move onto the next clause}
            \EndIf
        \EndFor
        \State $solved \gets satisfied \wedge solved$
        \If{$solved = false$} 
            \State \texttt{Move onto the next assignment}
        \EndIf
    \EndFor

    \If{$solved = True$}
        \State $vars \gets assignment$
        \State return $True$
    \EndIf
\EndFor

\State return $False$
\EndProcedure
\end{algorithmic}
\end{algorithm}


With that, that's all the optimization I can think of. In case you're
wondering, no this does not change the asymptotic runtime of the method,
however, it's easy to see why these changes would presumably make the
algorithm run faster as we're adding code that stops evaluation the moment
we've figured just enough out to know if it's worth continuing evaluation and
deferring extra calculation (recording the current assignment) until once we
know it's satisfiable.

As an exercise to the reader, I will ask you, how can you parallelize this
algorithm? And once you've figured that out, can you program it in a way
where once one thread finds a satisfactory assignment, it terminates all other
threads immediately instead of waiting for the others to finish?

\subsection{Generating Assignments}

How do we generate assignments? In the C++ implementation, we don't generate an assignment 
directly. Instead, we can recognize that a string of 1's (true) and 0's (false) is just a 
string of bits, which is just a datatype. We can use a large data, such as a \textit{long}
to hold 32 variable assignments. Since this is just a data type, if we had a for loop from 
$i \rightarrow 2^32$, or however many variables we have involved, our index variable can be used 
as an assignment itself! Using some clever bit shifting, we can either store the assignment into 
a bit array or index the bits of $i$ directly to simulate the behavior of accessing the assignment 
of a variable. You may think that 32 variables, or 64, depending on your compiler and architecture 
is very small, and while that may be true, you must consider how inefficient of an algorithm brute 
force is to begin with. Theoretically we \textit{could} make brute force numerically scale to any 
number of inputs if we wanted to, but the algorithm is so terrible that after 32, it would become 
basically unusable anyways other than a potential CPU hog.

\section{Brute Force - Results}

To see how well our current algorithm does, running some tests would be a
good idea. Why run tests when we have already proven that our algorithm
works ``perfectly?'' To see how quickly it runs of course.

To test its speed (and as we add more features, to ensure correctness is
kept), we use a database of problems from here:

\begin{center}
    \url{https://www.cs.ubc.ca/~hoos/SATLIB/benchm.html}
\end{center}

We gave our current implementation 1000 SAT instances (uf20-91), each
with 20 variables and 91 clauses each, all satisfiable. It took this baseline
brute force algorithm \textbf{2,119,301 milliseconds}, which is roughly 35.32
minutes. This result is actually from the first time I wrote a brute force 
algorithm. The brute force result you can find on GitHub, it actually took 
\textbf{1,770,851 milliseconds}, which is about 29.51 minutes. I have no idea what I 
did in the rewrite that shaved off 4 minutes. I won't even lie to you reader,
I am a busy undergraduate student, so I didn't even do the "best practice" for 
testing, which would be to quit all other processes and then benchmark it. However,
while I was testing, I was doing a Canvas assignment, so I was doing lots of 
background work.

Normally in science you have to run an experiment at least three times for
data accuracy. But you must understand that I don't want to give up an hour
and a half of my life to an algorithm that we're going to be improving upon
anyways.

Here are the results for all brute force algorithms, including the short-circuited 
and parallel versions:

\begin{FVerbatim}
On all 1000 equations in uf20-91:

Brute Force:               1,770,851ms (~29.51 minutes).
Short-Circuited (Basic):   49,874ms (~50 seconds).
Parallel:                  10,589ms (~10.5 seconds).
\end{FVerbatim}

For the parallel version, the way it works is that it spawns in some number of 
worker threads and they all share a flag that alerts them if any other thread has 
found a solution. If a thread has, each worker will finish whatever iteration it's 
on and terminate. The full code can be found on the GitHub.For the parallel version, the 
implementation isn't perfect. I believe that we don't prevent false sharing amongst cache 
lines (ie we get more cache swaps between threads leading to more cache misses), which could 
explain why on my M1 MacBook Air, we don't get close to a ~8 times speedup. Since this was for 
fun, I'm not going to perfect the parallelized version, that can be an exercise to the reader :\verb|)|